\section*{अध्याय \ref{ch:b1:elimination} --- \texorpdfstring{$\beta$}{बीटा}-विलोपन}

\begin{solution}{exo:b1:elimination:1}
2-ब्रोमोब्यूटेन: ब्यूट-1-ईन, \cip{E}- और \cip{Z}-ब्यूट-2-ईन; मुख्य
\cip{E}-ब्यूट-2-ईन। 2-ब्रोमो-2-मेथिलब्यूटेन: 2-मेथिलब्यूट-2-ईन (मुख्य,
त्रिप्रतिस्थापित) और 2-मेथिलब्यूट-1-ईन। ब्रोमोसाइक्लोहेक्सेन: केवल साइक्लोहेक्सीन।
\end{solution}

\begin{solution}{exo:b1:elimination:2}
E1: $v = k[\mathrm{RBr}]$; E2: $v = k[\mathrm{RBr}][\ce{EtO-}]$। कई एथॉक्साइड
सांद्रताओं पर ऐल्कीन बनने की \omterm{def:b1:rate-laws:initial-rate}{प्रारंभिक दर} मापिए: E2 के लिए वह
समानुपात में बढ़ती है, E1 के लिए नहीं बदलती।
\end{solution}

\begin{solution}{exo:b1:elimination:3}
C2--C3 की दिशा में, C2 पर Br, C3 के किसी भी हाइड्रोजन के प्रति या
C4 मेथिल समूह के प्रति प्रति-स्थिति में हो सकता है: जिन दो \omterm{def:b1:stereochemistry-in-depth:isomers}{संरूपणों} में कोई H, Br के प्रति प्रति-स्थिति में है, वे
E2 से अभिक्रिया कर सकते हैं (\cip{E} और \cip{Z} ऐल्कीन देते हुए); तीसरा
C3 की ओर विलोपन नहीं कर सकता।
\end{solution}

\begin{solution}{exo:b1:elimination:4}
SN2 (प्राथमिक, हाइड्रॉक्साइड, निम्न ताप); E2 (तृतीयक, प्रबल भारी-भरकम
क्षारक); SN2 (आयोडाइड, अच्छा \omterm{def:b1:electronic-effects:nucleophile}{नाभिकरागी} और \omterm{def:b1:predominant-reaction:strong-acid}{दुर्बल क्षारक}, \omterm{def:b1:intermolecular-forces-solvents:solvent-classes}{अप्रोटिक विलायक});
SN1 और E1, गरम करने के साथ E1 बढ़ता है।
\end{solution}

\begin{solution}{exo:b1:elimination:5}
C3 पर एक ही हाइड्रोजन है। हर \omterm{def:b1:stereochemistry-in-depth:enantiomers}{अप्रतिबिंबरूपी} में, उस H को C2 के Br के प्रति
प्रति-स्थिति में रखने से अन्य समूहों की स्थितियाँ निर्धारित हो जाती हैं: C2 का मेथिल
और C3 का एथिल एक \omterm{def:b1:stereochemistry-in-depth:enantiomers}{अप्रतिबिंबरूपी} में एक ही ओर हैं और दूसरे में
विपरीत ओर। एक \cip{E}- देता है, दूसरा
\cip{Z}-3-मेथिलपेंट-2-ईन: एक \omterm{def:b1:nucleophilic-substitution:stereospecific}{त्रिविमविशिष्ट अभिक्रिया}।
\end{solution}

\begin{solution}{exo:b1:elimination:6}
OH प्रोटॉनित होता है, \ce{R-OH2+}; जल निकलता है, जिससे तृतीयक
\omterm{def:b1:electronic-effects:intermediates}{कार्बधनायन} \ce{(CH3)2C+-CH2CH3} बनता है; एक क्षारक (जल, \ce{HSO4-}) पड़ोसी कार्बन से
एक प्रोटॉन हटाता है। मुख्य उत्पाद (ज़ैत्सेव):
2-मेथिलब्यूट-2-ईन।
\end{solution}

\begin{solution}{exo:b1:elimination:7}
टर्ट-ब्यूटिल समूह \omterm{def:b1:stereochemistry-in-depth:conformations}{कुर्सी} को उस रूप में बाँध देता है जिसमें वह स्वयं \omterm{def:b1:stereochemistry-in-depth:conformations}{विषुवतीय} है।
\textit{trans} समावयवी में तब ब्रोमीन भी \omterm{def:b1:stereochemistry-in-depth:conformations}{विषुवतीय} होता है, किसी C--H के प्रति
कभी प्रति-स्थिति में नहीं: E2 को अत्यंत प्रतिकूल पलटी हुई \omterm{def:b1:stereochemistry-in-depth:conformations}{कुर्सी} की प्रतीक्षा करनी पड़ती है।
\textit{cis} समावयवी में ब्रोमीन \omterm{def:b1:stereochemistry-in-depth:conformations}{अक्षीय} है, दो trans-द्विअक्षीय
हाइड्रोजनों के साथ: तीव्र E2।
\end{solution}

\begin{solution}{exo:b1:elimination:8}
छोटा एथॉक्साइड वह हाइड्रोजन हटाता है जो अधिक प्रतिस्थापित
ऐल्कीन देता है (ज़ैत्सेव)। भारी-भरकम टर्ट-ब्यूटॉक्साइड मेथिल समूह के खुले हाइड्रोजनों तक
आंतरिक \ce{CH2} के हाइड्रोजन की तुलना में अधिक आसानी से
पहुँचता है: कम प्रतिस्थापित ऐल्कीन।
\end{solution}

\begin{solution}{exo:b1:elimination:9}
ब्रोमोमेथेन में कोई $\beta$ कार्बन नहीं है। 1-ब्रोमो-2,2-डाइमेथिलप्रोपेन में
$\beta$ कार्बन चतुष्क है और उस पर कोई हाइड्रोजन नहीं है।
\end{solution}

\begin{solution}{exo:b1:elimination:10}
दोनों उत्पादों को \omterm{def:b1:electronic-effects:intermediates}{कार्बधनायन} चाहिए, जो मंद चरण में
$k_1[\mathrm{RBr}]$ दर से बनता है। फिर वह एथेनॉल द्वारा ग्रहण (\omterm{def:b1:rate-laws:order}{दर
स्थिरांक} $k_S$) और प्रोटॉन की हानि (\omterm{def:b1:rate-laws:order}{दर स्थिरांक} $k_E$) के बीच बँट जाता है, दोनों
विलायक में आभासी प्रथम कोटि के: ऐल्कीन का अंश
$k_E/(k_E + k_S)$ है, जो \omterm{def:b1:nucleophilic-substitution:substitution}{क्रियाधार} की सांद्रता से स्वतंत्र है।
\end{solution}

\begin{solution}{exo:b1:elimination:11}
\omterm{def:b1:stereochemistry-in-depth:enantiomers}{मेसो यौगिक} के उस \omterm{def:b1:stereochemistry-in-depth:isomers}{संरूपण} से आरंभ करके जिसमें दोनों Br प्रति-स्थिति में हैं
(तब दोनों फ़ेनिल भी प्रति-स्थिति में, दोनों H भी), C1 को $120^\circ$ घुमाइए ताकि
उसका H, C2 के Br के प्रति प्रति-स्थिति में आ जाए। तब दोनों फ़ेनिल समूह H--Br अक्ष की
एक ही ओर होते हैं: वे \textit{cis} पहुँचते हैं। C1 पर Br का क्रम
Ph से ऊपर है; C2 पर Ph का क्रम H से ऊपर; Br और C2 का फ़ेनिल \textit{trans} हैं:
\cip{E}-1-ब्रोमो-1,2-डाइफ़ेनिलएथीन।
\end{solution}

\begin{solution}{exo:b1:elimination:12}
केवल वह \omterm{def:b1:stereochemistry-in-depth:conformations}{कुर्सी} अभिक्रिया करती है जिसमें \omterm{def:b1:electronic-effects:nucleophile}{निर्गामी समूह} \omterm{def:b1:stereochemistry-in-depth:conformations}{अक्षीय} है: $v = k\,x\,[\mathrm{RY}][\mathrm{B}]$,
जहाँ $x$ उस \omterm{def:b1:stereochemistry-in-depth:conformations}{कुर्सी} का अंश है। उसके लिए आवश्यक हर \omterm{def:b1:stereochemistry-in-depth:conformations}{अक्षीय} ऐल्किल समूह
कई kJ/mol की कीमत लेता है (मेथिल के लिए \qty{5.7}{kJ/mol}), जो कमरे के ताप पर $x$ को
प्रति समूह लगभग दस से भाग देता है: ऐसे दो या तीन समूहों के साथ $x$
$10^{-2}$ या $10^{-3}$ तक गिर जाता है और अभिक्रिया उतनी ही धीमी हो जाती है।
\end{solution}

\begin{solution}{pb:b1:elimination:1}
\textbf{1.} ऐसी \omterm{def:b1:stereochemistry-in-depth:conformations}{कुर्सी} जिसमें Cl, आइसोप्रोपिल और मेथिल तीनों \omterm{def:b1:stereochemistry-in-depth:conformations}{विषुवतीय} हैं।
\textbf{2.} आइसोप्रोपिल और मेथिल \omterm{def:b1:stereochemistry-in-depth:conformations}{विषुवतीय}, Cl \omterm{def:b1:stereochemistry-in-depth:conformations}{अक्षीय}।
\textbf{3.} \omterm{def:b1:stereochemistry-in-depth:enantiomers}{अप्रतिबिंबरूपी}: वे केवल C1 पर भिन्न हैं।
\textbf{4.} \omterm{def:b1:electronic-effects:nucleophile}{निर्गामी समूह} और $\beta$ हाइड्रोजन trans-द्विअक्षीय।
\textbf{5.} केवल पलटी हुई \omterm{def:b1:stereochemistry-in-depth:conformations}{कुर्सी} में, जिसमें Cl, आइसोप्रोपिल और मेथिल तीनों
\omterm{def:b1:stereochemistry-in-depth:conformations}{अक्षीय} हैं।
\textbf{6.} तीनों समूह \omterm{def:b1:stereochemistry-in-depth:conformations}{अक्षीय}: अकेला मेथिल लगभग
\qty{5.7}{kJ/mol} की कीमत लेता है, बड़ा आइसोप्रोपिल समूह अधिक, क्लोरीन थोड़ी:
कुल मिलाकर लगभग \qty{13}{kJ/mol}, इसलिए अंश लगभग
$\exp(-13000/2479) = \num{5e-3}$, वही मान जो समस्या लेती है।
\textbf{7.} C2 (एक H, वह आइसोप्रोपिल समूह वहन करता है) और C6 (दो H)।
\textbf{8.} C1 पर Cl \omterm{def:b1:stereochemistry-in-depth:conformations}{अक्षीय} होने पर, C2 का \omterm{def:b1:stereochemistry-in-depth:conformations}{अक्षीय} H (आइसोप्रोपिल
\omterm{def:b1:stereochemistry-in-depth:conformations}{विषुवतीय} होने से) और C6 का \omterm{def:b1:stereochemistry-in-depth:conformations}{अक्षीय} H: दो।
\textbf{9.} त्रिअक्षीय \omterm{def:b1:stereochemistry-in-depth:conformations}{कुर्सी} में आइसोप्रोपिल समूह C2 की अक्षीय
स्थिति घेरता है, इसलिए उसका H \omterm{def:b1:stereochemistry-in-depth:conformations}{विषुवतीय} है: प्रति-स्थिति में नहीं। C6 पर एक \omterm{def:b1:stereochemistry-in-depth:conformations}{अक्षीय} H है: प्रति-स्थिति में।
\textbf{10.} निओमेंथिल क्लोराइड दो, मेंथिल क्लोराइड एक।
\textbf{11.} C1--C6 की दिशा में: C1 पर Cl (\omterm{def:b1:stereochemistry-in-depth:conformations}{अक्षीय}, ऊपर), C6 पर \omterm{def:b1:stereochemistry-in-depth:conformations}{अक्षीय} H (नीचे) के
विपरीत; वलय के आबंध और \omterm{def:b1:stereochemistry-in-depth:conformations}{विषुवतीय} H बीच में सांतरित।
\textbf{12.} \omterm{def:b1:stereochemistry-in-depth:conformations}{कुर्सी} पर कोई C--H आबंध C--Cl आबंध के प्रति \omterm{def:b1:elimination:periplanar}{सम-परिसमतलीय} नहीं है;
उसे बलपूर्वक लाने का अर्थ होगा एक ग्रसित, तनावयुक्त वलय।
\textbf{13.} C2 का H हटाने पर: 4-मेथिल-1-(प्रोपेन-2-इल)साइक्लोहेक्स-1-ईन
(त्रिप्रतिस्थापित); C6 का H हटाने पर: 3-मेथिल-6-(प्रोपेन-2-इल)साइक्लोहेक्स-1-ईन
(द्विप्रतिस्थापित)।
\textbf{14.} त्रिप्रतिस्थापित ऐल्कीन, जो अधिक प्रतिस्थापित और स्थायी है।
\textbf{15.} दोनों ऐल्कीन, \qty{78}{\percent} त्रिप्रतिस्थापित: संगत।
\textbf{16.} केवल द्विप्रतिस्थापित ऐल्कीन: ज़ैत्सेव के
पूर्वानुमान के विपरीत, जिसे एकमात्र उपलब्ध प्रति हाइड्रोजन थोपता है।
\textbf{17.} दोनों त्रिविमविशिष्ट हैं (प्रति अपेक्षा); निओमेंथिल क्लोराइड
स्थानवरणात्मक रूप से अभिक्रिया करता है (78/22); मेंथिल क्लोराइड एक अकेला
\omterm{def:b1:stereochemistry-in-depth:isomers}{संरचनात्मक समावयवी} देता है।
\textbf{18.} Cl वहन करने वाला कार्बन द्वितीयक है और एक आइसोप्रोपिल
समूह से सटा है: SN2 के लिए भीड़, जबकि एथॉक्साइड \omterm{def:b1:predominant-reaction:strong-acid}{प्रबल क्षारक} है।
\textbf{19.} $v \propto x\,n_{\mathrm{H}}$: क्रियाशील \omterm{def:b1:stereochemistry-in-depth:conformations}{कुर्सी} का अंश गुणा
प्रति हाइड्रोजनों की संख्या।
\textbf{20.} $0.95 \times 2 = 1.90$।
\textbf{21.} $\num{5.0e-3} \times 1 = \num{5.0e-3}$।
\textbf{22.} \textbf{$k(\text{निओमेंथिल})/k(\text{मेंथिल}) = 1.90/\num{5.0e-3}
= 380$}।
\end{solution}
